\documentclass[11pt,reqno]{amsart}

%==============================================================================
%  Two-sided bounds on the essential minimum of the Zhang--Zagier height
%
%  Compile with pdflatex (or xelatex). No non-standard packages required.
%==============================================================================

\usepackage[utf8]{inputenc}
\usepackage[T1]{fontenc}
\usepackage{amsmath,amssymb,amsthm,mathtools}
\usepackage[colorlinks=true,linkcolor=blue,citecolor=blue,urlcolor=blue]{hyperref}
\usepackage{url}
\usepackage{booktabs}

%------------------------------------------------------------------ theorems
\theoremstyle{plain}
\newtheorem{theorem}{Theorem}
\newtheorem{proposition}[theorem]{Proposition}
\newtheorem{lemma}[theorem]{Lemma}
\newtheorem{corollary}[theorem]{Corollary}
\theoremstyle{definition}
\newtheorem{remark}[theorem]{Remark}

%------------------------------------------------------------------ macros
\newcommand{\Q}{\mathbb{Q}}
\newcommand{\Z}{\mathbb{Z}}
\newcommand{\R}{\mathbb{R}}
\newcommand{\C}{\mathbb{C}}
\newcommand{\Qbar}{\overline{\Q}}
\newcommand{\hZ}{h_Z}
\newcommand{\Cee}{C_{82}}
\newcommand{\disc}{\operatorname{disc}}
\newcommand{\Res}{\operatorname{Res}}
\newcommand{\Li}{\operatorname{Li}}
\newcommand{\Cl}{\operatorname{Cl}}
\newcommand{\Ihat}{\widehat{I}}
\newcommand{\dnu}{\,d\nu}
\DeclareMathOperator{\Real}{Re}

\newcommand{\repo}{\url{https://github.com/tempcollab/optimizationproblems/tree/main/constants/82a}}
\newcommand{\repofile}[1]{\href{https://github.com/tempcollab/optimizationproblems/blob/main/constants/82a/certificate/#1}{\texttt{#1}}}

\begin{document}

\title[Two-sided bounds on the Zhang--Zagier essential minimum]
{Two-sided bounds on the essential minimum\\ of the Zhang--Zagier height}

\author{Adib Hasan}
\address{SignalPilot AI}
\email{adib.hasan8@gmail.com}

\author{Akashnil Dutta}
\address{Incept Labs}

\author{Tarik Moon}
\address{SignalPilot AI}

\date{\today}

\subjclass[2020]{11R06, 11G50, 11Y40 (primary); 31A15 (secondary)}
\keywords{Zhang--Zagier height, essential minimum, Mahler measure, logarithmic
energy, discriminant, auxiliary functions, interval arithmetic}

\begin{abstract}
Let $\Cee$ denote the essential minimum of the Zhang--Zagier height
$\hZ(\alpha)=h(\alpha)+h(1-\alpha)$ on $\Qbar$. We prove
\[
  0.2524001332 \;\le\; \Cee \;\le\; 0.2540419719,
\]
improving both classical records: the lower bound $0.2487458$ of
Flammang (2018) and the upper bound $0.25443677$ of Doche (2001). The lower
bound is obtained by adjoining to Smyth's auxiliary-function method a new dual
constraint coming from the Orloski--Sardari--Smith logarithmic-energy /
discriminant inequality; this is the first transfer of that constraint to the
Zhang--Zagier setting and the first improvement of the lower bound since 2018.
The upper bound enriches Doche's perturbed-polynomial limit-point family with
two additional admissible perturbing blocks. Every numerical inequality is
established by an outward-rounded interval-arithmetic certificate that the reader
can re-run; the certificates are available at \repo.
\end{abstract}

\maketitle

%==============================================================================
\section{Introduction}
%==============================================================================

For $\alpha\in\Qbar$ let $h(\alpha)$ be the absolute logarithmic Weil height. The
\emph{Zhang--Zagier height} is
\[
  \hZ(\alpha)\;=\;h(\alpha)+h(1-\alpha),
\]
a height that simultaneously measures the arithmetic complexity of $\alpha$ near
$0$ and near $1$. Zhang \cite{Zha95} and Zagier \cite{Zag93} initiated its study;
Zagier proved $\hZ(\alpha)\ge\tfrac12\log\frac{1+\sqrt5}{2}$ for
$\alpha\notin\{0,1,\tfrac{1\pm\sqrt{-3}}{2}\}$. The quantity of interest here is the
\emph{essential minimum}
\begin{equation}\label{eq:def}
  \Cee\;=\;\sup\Bigl\{\,H\in\R:\ \{\alpha\in\Qbar:\hZ(\alpha)\le H\}\ \text{is finite}\,\Bigr\},
\end{equation}
the largest height threshold below which only finitely many algebraic numbers
lie. Its exact value is unknown; the problem is to narrow the interval
$[\,\text{lower},\text{upper}\,]$ bracketing it.

The two sides are approached by dual methods. Lower bounds come from Smyth's
auxiliary-function method \cite{S80}, developed for $\hZ$ by Doche \cite{Doc01a}
and Flammang \cite{F18}, whose bound $\Cee\ge\log 1.282416=0.2487458$ has stood
since 2018. Upper bounds come from explicit limit-point constructions; Doche's
perturbed-polynomial family \cite{Doc01b} gives $\Cee\le\log 1.289735=0.25443677$.
Burgos Gil, Menares, Qu and Sombra \cite{BMQS} recently proved that these two
methods are linear-programming duals of one another, with \emph{no duality gap},
so the residual interval reflects only how far the two finite optimizations have
been pushed.

Our main result improves both records simultaneously.

\begin{theorem}\label{thm:main}
The essential minimum of the Zhang--Zagier height satisfies
\[
  0.2524001332\;\le\;\Cee\;\le\;0.2540419719 .
\]
\end{theorem}

The lower bound (Theorem~\ref{thm:lower}) strictly exceeds the Flammang record by
$3.654\times10^{-3}$; the upper bound (Theorem~\ref{thm:upper}) strictly improves
the Doche record by $3.42\times10^{-4}$. The resulting interval has width
$\approx1.64\times10^{-3}$, down from $\approx5.7\times10^{-3}$ before this work.

The lower bound is the principal contribution. It adjoins to Flammang's
auxiliary function a constraint that her method does not carry: the
\emph{logarithmic-energy / discriminant inequality} of Orloski, Sardari and Smith
\cite{OSS}, the same constraint that produced the largest advance on the
Schur--Siegel--Smyth trace problem in forty years. Transferred to the
Zhang--Zagier conjugate measure, it is an independent dual column that can only
raise the bound. To our knowledge this is its first use outside the trace
problem.

\medskip
\noindent\textbf{On rigor and reproducibility.}
Every numerical inequality below is certified by \emph{outward-rounded interval
arithmetic}: each floating-point operation in the verifying program is replaced
by a guaranteed enclosure (directed rounding via \texttt{numpy.nextafter}), so no
assertion depends on an unrounded float. The programs, the frozen data file, and
a one-line reproduction command for each theorem are collected in
Section~\ref{sec:repro} and hosted at \repo.

%==============================================================================
\section{Setup and notation}\label{sec:setup}
%==============================================================================

If $\alpha$ is an algebraic integer of degree $d$ with conjugates
$\alpha_1,\dots,\alpha_d$ and minimal polynomial $P\in\Z[z]$, then
\[
  \hZ(\alpha)\;=\;\tfrac1d\log Z(\alpha),\qquad
  Z(\alpha)\;=\;M(\alpha)\,M(1-\alpha),
\]
with $M$ the Mahler measure; we also write $\zeta(\alpha)=Z(\alpha)^{1/d}$, so
$\hZ(\alpha)=\log\zeta(\alpha)$.

We work in the symmetric coordinate
\[
  w\;=\;z(1-z),
\]
invariant under $z\mapsto 1-z$. On the unit circle $z=e^{it}$ we set
$w(t)=e^{it}-e^{2it}$, $t\in[0,\pi]$. The Zhang--Zagier Green function is
\[
  g(z)\;=\;\log^+|z|+\log^+|1-z|,\qquad \log^+x=\log\max(1,x),
\]
and on $|z|=1$ one has $g|_{\{|z|=1\}}=\log^+|w(t)|$.

For an algebraic integer $\alpha\notin\{0,1\}$ of degree $d$ the \emph{conjugate
measure} is $\nu=\frac1d\sum_{i=1}^d\delta_{\alpha_i}$ (in the $z$-coordinate);
since $P\in\Z[z]$, $\nu$ is invariant under complex conjugation.

We use the duality of \cite[Thm.~6.5]{BMQS}: with $g$ as above,
\[
  \underbrace{D(g)}_{\text{auxiliary-function (lower)}}
  \;=\;\Cee\;=\;
  \underbrace{P(g)}_{\text{measure (upper)}},
\]
strong duality with no gap. Here $D(g)$ is the supremum over $c_j\ge0$,
$Q_j\in\Z[x]\setminus\{0\}$ of $\inf_z\bigl(g(z)-\sum_jc_j\log|Q_j(z)|\bigr)$, and
$P(g)$ is the infimum of $\int g\,d\mu$ over conjugation-invariant probability
measures $\mu$ with $\int\log|Q|\,d\mu\ge0$ for all $Q\in\Z[x]\setminus\{0\}$.

\medskip
\noindent\textbf{The records we improve.}
The best previously \emph{verified} bounds are the Flammang lower bound
$0.2487458=\log 1.282416$ \cite{F18} and the Doche upper bound
$0.25443677=\log1.289735$ \cite{Doc01b}. (The structural paper \cite{BMQS} quotes
the weaker pair $0.248247\le\Cee\le0.254437$ and does not cite Flammang; the
genuine bar is the Flammang/Doche pair.)

%==============================================================================
\section{The auxiliary-function method and the integrality columns}\label{sec:aux}
%==============================================================================

We recall Smyth's method as used by Flammang, isolating the load-bearing
integrality step, because the new energy column of Section~\ref{sec:energy} rests
on the same step.

Fix weights $c_j\ge0$ and integer polynomials $Q_j\in\Z[w]$, $1\le j\le J$, a real
$\lambda_0\ge0$, and a nonnegative \emph{energy column} $E(z)\ge0$ to be defined in
Section~\ref{sec:energy} (set $\lambda_0=0$ for the classical method). Consider
\begin{equation}\label{eq:aux}
  f(z)\;=\;g(z)\;-\;\sum_{j=1}^{J}c_j\log|Q_j(w)|\;-\;\lambda_0\,E(z),
  \qquad w=z(1-z).
\end{equation}

\begin{proposition}[Drop-out]\label{prop:dropout}
Let $\alpha$ be an algebraic integer of degree $d$ whose minimal polynomial $P$
divides none of the polynomials $Q_j(z(1-z))$, and suppose the energy column
satisfies $\int E\dnu\ge0$. If $f(z)\ge m$ for all $z$ with $|z|=1$, then
$\hZ(\alpha)\ge m$.
\end{proposition}

\begin{proof}
Average \eqref{eq:aux} against the conjugate measure $\nu$. By definition
$\int g\dnu=\frac1d\log Z(\alpha)=\hZ(\alpha)$. For each column,
\[
  \int\log|Q_j(w)|\dnu
  =\tfrac1d\log\Bigl|\textstyle\prod_{i=1}^d Q_j\bigl(\alpha_i(1-\alpha_i)\bigr)\Bigr|
  =\tfrac1d\log\bigl|\Res\bigl(P(z),Q_j(z(1-z))\bigr)\bigr|.
\]
The resultant is the determinant of an integer Sylvester matrix, hence an integer;
it is nonzero precisely because $P\nmid Q_j(z(1-z))$. Thus $|\Res|\ge1$ and
\begin{equation}\label{eq:colpos}
  \int\log|Q_j(w)|\dnu\;\ge\;0 .
\end{equation}
With $c_j\ge0$ these terms are subtracted with nonnegative weight, and likewise
$\lambda_0\ge0$, $\int E\dnu\ge0$. Hence
\[
  \hZ(\alpha)=\int g\dnu\;\ge\;\int f\dnu\;\ge\;m. \qedhere
\]
\end{proof}

\paragraph{The finite exception set}
Proposition~\ref{prop:dropout} fails only for the finitely many $\alpha$ that are
roots of some fixed nonzero $Q_j(z(1-z))$ (or, for the energy column, of the
discriminant exception of Section~\ref{sec:energy}). Let $E_{\mathrm{exc}}$ be the
finite union of these root sets; then $\hZ(\alpha)\ge m$ for all algebraic
integers $\alpha\notin E_{\mathrm{exc}}$.

\paragraph{A finite exception set does not lower \texorpdfstring{$\Cee$}{C82}}
If $\hZ(\alpha)\ge m$ for all $\alpha\notin E_{\mathrm{exc}}$ with
$E_{\mathrm{exc}}$ finite, then for every $H<m$ the sublevel set
$\{\alpha:\hZ(\alpha)\le H\}\subseteq E_{\mathrm{exc}}$ is finite, so every $H<m$
is admissible in \eqref{eq:def} and $\Cee\ge m$. The restriction to algebraic
integers is harmless: for non-integers the leading coefficient only raises $h$.

\paragraph{Reduction to the unit circle}
$f$ is harmonic off small disks around the roots of the $Q_j(z(1-z))$, together
with the single term $-2\lambda_0 U_{\mu_0}$ which is the logarithmic potential of
a measure supported on $|z|=1$ (Section~\ref{sec:energy}), hence harmonic off
$|z|=1$ and superharmonic across it. By the minimum principle $f$ attains its
minimum over the lens $\{|z|\le1\}\cap\{|1-z|\le1\}$ on the boundary; the symmetry
$f(z)=f(1-z)$ (since $z\mapsto1-z$ fixes $w$ and $E$ is built symmetrically) then
reduces the problem to minimizing
\begin{equation}\label{eq:f-circle}
  f|_{\{|z|=1\}}(t)
  =\log^+|w(t)|-\sum_jc_j\log|Q_j(w(t))|-\lambda_0\bigl(2U_{\mu_0}(e^{it})-\Ihat\bigr),
  \quad t\in[0,\pi].
\end{equation}

With $\lambda_0=0$ and Flammang's $J=24$ Table~1 columns, the rigorous minimum of
\eqref{eq:f-circle} is $\ge0.2487458$; we re-certified this from scratch
(\repofile{verify\_vec.py}, \repofile{verify\_flammang.py}), independently
re-deriving \eqref{eq:colpos}. This is the base to which we add the energy column.

%==============================================================================
\section{The logarithmic-energy / discriminant column}\label{sec:energy}
%==============================================================================

\subsection{The hard energy constraint}
For the conjugate measure $\nu$ of a minimal polynomial $P$ of degree $d$,
\begin{equation}\label{eq:Inu}
  I(\nu)\;:=\;\iint\log|z_1-z_2|\,d\nu(z_1)\,d\nu(z_2)
  \;=\;\frac1{d^2}\log\bigl|\disc(P)\bigr|.
\end{equation}
Since $\disc(P)=\prod_{i<j}(\alpha_i-\alpha_j)^2$ is a nonzero integer for every
separable $P\in\Z[z]$ off a finite exception set, $|\disc(P)|\ge1$ and
\begin{equation}\label{eq:Ipos}
  I(\nu)\;\ge\;0 .
\end{equation}
This is precisely the integrality clause of Proposition~\ref{prop:dropout}, now
applied to the bivariate integer polynomial $Q(z_1,z_2)=z_1-z_2$
\cite{OSS,OT23}. It is independent of the univariate columns
\eqref{eq:colpos}: the energy couples two copies of $\nu$, and no finite family of
univariate columns reproduces it. Flammang's dual does not carry it.

\subsection{Linearization: the supergradient outer cut}
Constraint \eqref{eq:Ipos} is quadratic in $\nu$, so it cannot enter a linear
auxiliary function directly. Following \cite[eq.~(8)]{OSS} we linearize it. The
energy $I(\cdot)$ is concave on finite-energy measures of fixed mass because the
logarithmic kernel is negative-definite (Ahlfors; \cite[Lem.~2.1, eq.~(6)]{OSS}):
for any two finite-energy probability measures $\nu,\mu_0$,
\begin{equation}\label{eq:nsd}
  I(\nu)+I(\mu_0)-2I(\nu,\mu_0)\;\le\;0,\qquad
  I(\nu,\mu_0):=\iint\log|z_1-z_2|\,d\nu(z_1)\,d\mu_0(z_2).
\end{equation}
With the logarithmic potential $U_{\mu_0}(z):=\int\log|z-z'|\,d\mu_0(z')$, so that
$I(\nu,\mu_0)=\int U_{\mu_0}\dnu$, rearranging \eqref{eq:nsd} gives
\begin{equation}\label{eq:cut}
  \int\bigl(2U_{\mu_0}(z)-I(\mu_0)\bigr)\dnu
  \;=\;2I(\nu,\mu_0)-I(\mu_0)\;\ge\;I(\nu)\;\ge\;0 .
\end{equation}
Let $\Ihat$ be any rigorous lower bound $\Ihat\le I(\mu_0)$. Replacing $I(\mu_0)$
by $\Ihat$ only adds the nonnegative quantity $I(\mu_0)-\Ihat$ to the left side of
\eqref{eq:cut}, so the \emph{energy column}
\begin{equation}\label{eq:Ecol}
  E(z):=2U_{\mu_0}(z)-\Ihat\qquad\text{satisfies}\qquad \int E\dnu\;\ge\;0 .
\end{equation}
Two points secure rigor:
\begin{itemize}
\item \emph{Validity rests on the continuous inequality \eqref{eq:nsd}, not on any
  discretization.} It holds whenever $\nu$ and $\mu_0$ both have finite energy:
  $\nu$ does off the discriminant exception set (where $I(\nu)<\infty$ by
  \eqref{eq:Inu}), and $\mu_0$ is chosen diffuse (Section~\ref{subsec:mu0}). An
  atomic $\mu_0$ would have $I(\mu_0)=-\infty$ and render \eqref{eq:cut} vacuous;
  this failure mode is avoided.
\item \emph{The cut can only raise the bound, never overshoot $\Cee$.} Because
  \eqref{eq:Ecol} holds for the conjugate measure of every admissible $\alpha$,
  adding $-\lambda_0 E$ with $\lambda_0\ge0$ to \eqref{eq:aux} preserves
  Proposition~\ref{prop:dropout}; as an outer linearization it gives
  $m_0\le m_{\mathrm{cut}}\le m_{\mathrm{true}}\le\Cee$.
\end{itemize}

\subsection{The diffuse reference measure}\label{subsec:mu0}
We take $\mu_0$ to be a histogram of uniform arcs on $|z|=1$: a piecewise-constant
density on $K=15$ arcs of common $t$-width $L=0.057119866$, with centers
$\{c_k\}$ and masses $\{\rho_k\}$, $\rho_k\ge0$, $\sum_k\rho_k=1$, symmetric under
$z\mapsto\bar z$. Being a bounded density it has finite energy, so
\eqref{eq:nsd}--\eqref{eq:Ecol} apply.

For uniform arcs $I(\mu_0)$ is exactly summable. Writing
\[
  M(c_a,c_b)=\tfrac12\bigl(B_1(c_a-c_b,L)+B_1(c_a+c_b,L)\bigr),\quad
  B_1(d,L)=\frac{F_2(|d|-L)+F_2(|d|+L)-2F_2(|d|)}{L^2},
\]
\[
  F_2(x)=\Real\Li_3(e^{ix})-\zeta(3),
\]
one has $I(\mu_0)=\sum_{a,b}\rho_a\rho_b\,M(c_a,c_b)$. (The derivation uses
$F_2''(x)=\log\bigl(2\sin\tfrac x2\bigr)$, the antiderivative for the arc-average
of $\log|e^{is}-e^{is'}|$; the singular self-block $d=0$ is the finite limit of an
integrable logarithmic singularity.) Evaluating in high precision and rounding
\emph{down} to a float yields the rigorous downward enclosure
\begin{equation}\label{eq:Ihat}
  \Ihat\;\le\;I(\mu_0) .
\end{equation}
For the frozen $\mu_0$ used below, the certificate ships $\Ihat=-0.2111616260$,
while an independent reference value gives $I(\mu_0)=-0.2111600446$, so
\eqref{eq:Ihat} holds with margin $1.58\times10^{-6}$. A numerical witness of
\eqref{eq:nsd} is also recorded: the exact-block conjugate-symmetric energy matrix
on a $240$-node grid has top eigenvalue $+4.0\times10^{-16}$ on the mass-zero
subspace (negative-definite to machine precision). This is a check; the proof is
the continuous \eqref{eq:nsd}.

\subsection{Freezing the weights}
The weights $(c_j,\lambda_0)$ are obtained once by solving the cut linear
program---Flammang's no-energy program with the single extra row \eqref{eq:Ecol}
appended---and reading its dual multipliers. The linear program is only a search:
the validity of the lower bound depends on $c_j\ge0$, $\lambda_0\ge0$ and
$\Ihat\le I(\mu_0)$, all checked directly, not on any optimality. For the frozen
data: $\lambda_0=0.040126675>0$, $\sum_jc_j=0.43728719$ with all $c_j\ge0$, and a
conjectural program value $m_{\mathrm{cut}}=0.2526110$. The frozen objects are
recorded in \repofile{frozen\_energy.npz} and produced/validated by
\repofile{freeze\_energy.py}.

%==============================================================================
\section{The lower bound}\label{sec:lowerproof}
%==============================================================================

\begin{theorem}\label{thm:lower}
$\Cee\ge 0.2524001332$, strictly exceeding the Flammang record $0.2487458$ by
$3.654\times10^{-3}$.
\end{theorem}

By Sections~\ref{sec:aux}--\ref{sec:energy}, Theorem~\ref{thm:lower} reduces to
certifying, rigorously over the whole continuum $t\in[0,\pi]$, that
\begin{equation}\label{eq:lowertarget}
  f(t)=\log^+|w(t)|-\sum_{j=1}^{24}c_j\log|Q_j(w(t))|
  -\lambda_0\bigl(2U_{\mu_0}(e^{it})-\Ihat\bigr)\;\ge\;0.2524001332 .
\end{equation}
With Proposition~\ref{prop:dropout} and the finite-exception argument, this gives
$\Cee\ge 0.2524001332$.

\begin{proof}
The program \repofile{verify\_vec\_energy.py} certifies \eqref{eq:lowertarget} by
an outward-rounded interval branch-and-bound on $[0,\pi]$. Each cell $[a,b]$ is
assigned a guaranteed lower bound $L([a,b])\le\min_{t\in[a,b]}f(t)$; cells with
$L\ge\text{target}$ are certified, the rest bisected, and the run succeeds only
when the frontier is emptied with no cell left at maximum depth. All arithmetic
uses directed rounding (\texttt{numpy.nextafter}), so $L$ is a true lower bound.

\emph{The $\log^+|w|$ and $\log|Q_j|$ terms} are enclosed by a second-order
Taylor (mean-value) model: per factor, $|Q_j(w(t))|^2$ is bounded by its midpoint
value and derivative (exact, outward-rounded) plus a curvature enclosure over the
cell. Only an \emph{upper} bound on $|Q_j|^2$ is needed, so $-c_j\log|Q_j|$ is
bounded below; this stays finite even at the zeros of $Q_j$, where the subtracted
term $\to+\infty$ and only helps.

\emph{The energy term} $-\lambda_0(2U_{\mu_0}-\Ihat)$ requires an \emph{upper}
bound on $U_{\mu_0}(e^{it})$ over each cell. For the histogram $\mu_0$,
\[
  U_{\mu_0}(e^{it})=\sum_k\rho_k\cdot\tfrac12\bigl(P_k(t)+\widetilde P_k(t)\bigr),
  \quad
  P_k(t)=\frac{\Cl_2(t-c_k-\tfrac L2)-\Cl_2(t-c_k+\tfrac L2)}{L},
\]
with $\Cl_2(x)=\operatorname{Im}\Li_2(e^{ix})$ the Clausen function and
$\widetilde P_k$ the conjugate-arc term (center $-c_k$); this is the exact
arc-average of $\log|e^{it}-e^{is}|$ over the support arc. A rigorous upper bound
on $P_k$ follows from a rigorous interval enclosure of $\Cl_2$, monotone between
its extrema $\pm\Cl_2(\pi/3)=\pm1.0149416$, evaluated from
\texttt{scipy.special.spence} and padded outward. Over-estimating $U_{\mu_0}$
over-subtracts, i.e.\ \emph{lowers} the certified bound---the safe direction. Near
a support node the subtracted $-2\lambda_0U_{\mu_0}\to-\infty$, so those cells
carry the largest $f$ and simply bisect; nothing auto-certifies.

Running the certificate yields
\begin{quote}\small\ttfamily
[OK] CERTIFIED  min\_t f(t) >= 0.2524001332\\
\quad cells=7858, max depth=4, time=4.1s, frontier resolved (0 unresolved)
\end{quote}
with binding cell at $t\approx0.489$. By \eqref{eq:lowertarget},
$\Cee\ge0.2524001332$.
\end{proof}

\paragraph{Soundness audits}
All reproduced (Section~\ref{sec:repro}):
(i) the interval lower bound never exceeds the true $f$---\texttt{selftest}
checks it against the exact $f$ computed two independent ways
(\texttt{scipy.special.spence} and \texttt{mpmath.polylog}) on $620$ generic plus
near-node cells, with $0/620+0/80$ violations and the two evaluations agreeing to
$1.3\times10^{-15}$;
(ii) the enclosure \eqref{eq:Ihat} and the witness for \eqref{eq:nsd} are
reproduced by \texttt{freeze\_energy.py check} on the committed measure
($\Ihat=-0.2111616260\le I(\mu_0)_{\mathrm{ref}}=-0.2111600446$; top eigenvalue
$+4.0\times10^{-16}$);
(iii) targets above the true ceiling ($\approx0.252555$) correctly fail to
resolve, so the certificate does not rubber-stamp---the certified value carries
$\approx1.5\times10^{-3}$ of margin below that ceiling.

\begin{remark}
Theorem~\ref{thm:lower} is, to our knowledge, the first rigorous transfer of the
Orloski--Sardari--Smith logarithmic-energy / discriminant constraint to the
Zhang--Zagier essential minimum, and the first improvement of its lower bound
since \cite{F18}. The energy column is independent of Flammang's univariate
columns and can only raise the bound.
\end{remark}

%==============================================================================
\section{The upper bound}\label{sec:upperproof}
%==============================================================================

\begin{theorem}\label{thm:upper}
$\Cee\le 0.2540419719$, strictly improving the Doche record $0.25443677$ by
$3.42\times10^{-4}$.
\end{theorem}

The upper bound uses the measure side $P(g)$ \cite[Thm.~6.5]{BMQS}: an explicit
conjugation-invariant probability measure $\mu$ with $\int\log|Q|\,d\mu\ge0$ for
all $Q\in\Z[x]\setminus\{0\}$ gives $\Cee\le\int g\,d\mu$. Doche realizes such
measures as asymptotic Galois-conjugate distributions of a perturbed-polynomial
limit-point family \cite{Doc01a,Doc01b}.

\subsection{The limit-point family}
Work in $X=z(1-z)$. Fix Doche's five base polynomials
$\{P_1,P_2,P_4,P_6,P_8\}\subset\Z[X]$ and the distinguished block
$Q_1Q_2\in\Z[X]$ of degree $56$ (the exact \cite{Doc01b} record family, with
$Q_1,Q_2$ its two degree-$28$ factors). For an exponent vector $q\in\Q_+^5$ and
free perturbing blocks $Q_k\in\Z[X]$ with exponents $q_k\ge0$, the limit measure
value is
\begin{equation}\label{eq:logh}
  \log h(q,\{q_k\})=\frac1{2\pi D}\int_0^{2\pi}G(t)\,dt,\qquad
  G(t)=\max\bigl(A(t),B(t)\bigr),
\end{equation}
\[
  A(t)=\sum_i q_i\log|P_i(w(t))|,\quad
  B(t)=\log|Q_1Q_2(w(t))|+\sum_k q_k\log|Q_k(w(t))|,
\]
with the normalization
\begin{equation}\label{eq:Dformula}
  D=\max\Bigl(\textstyle\sum_i q_i\deg P_i,\ \ 56+\sum_k q_k\deg Q_k\Bigr).
\end{equation}
The form $G=\max(A,B)$ is the inner Jensen integral of Doche's double integral:
$\int_0^1\log|A-e^{2\pi it}B|\,dt=\log\max(|A|,|B|)$. Equation \eqref{eq:Dformula}
is the \cite[\S4]{Doc01a} $D$-formula.

\begin{lemma}[{\cite[Lemmas 2--5]{Doc01a}}]\label{lem:doche}
For any admissible family (Section~\ref{subsec:admiss}), $\log h(q,\{q_k\})$ is a
limit point of the spectrum $V=\{\hZ(\alpha)\}$: there is a sequence of distinct
irreducible algebraic numbers whose Zhang--Zagier heights converge to it.
Consequently
\begin{equation}\label{eq:upperbound}
  \Cee\;\le\;\log h(q,\{q_k\})\qquad\text{for every admissible }(q,\{q_k\}).
\end{equation}
\end{lemma}

No optimality of $q$ is required; we exploit this by enriching the dictionary with
new admissible blocks rather than re-tuning $q$.

\subsection{Admissibility}\label{subsec:admiss}
A perturbing block is admissible---so \eqref{eq:upperbound} holds---iff
\cite[Lem.~5 and condition (4)]{Doc01a}:
\begin{enumerate}
\item $\deg(Q_1Q_2)=56>0$;
\item each $Q_k$ satisfies $Q_k(0)=Q_k(1)=1$ (neither $X$ nor $1-X$ divides it);
\item each $Q_k$ is squarefree (irreducible suffices);
\item $\gcd(Q_k,\cdot)=1$ against every kept factor $P_1,P_2,P_4,P_6,P_8,Q_1,Q_2$
  and against every other active block;
\item $\prod P^n\ne\prod Q^n$, which holds by unique factorization once all factors
  are non-constant and pairwise coprime.
\end{enumerate}

\subsection{A fresh source of admissible blocks}
The structural observation is that an admissible block may be \emph{any} integer
polynomial in $X=z(1-z)$ meeting Section~\ref{subsec:admiss}; it need not come from
Doche's calibration set. We draw two blocks from Flammang's \cite{F18} Table~1
(their roles on the lower side are irrelevant here):
\begin{align*}
  Q_5\ (\deg 12):\ &X^{12}-3X^{11}+8X^{10}-18X^9+36X^8-62X^7+97X^6-123X^5\\
    &\quad{}+114X^4-73X^3+31X^2-8X+1,\\
  Q_6\ (\deg 16):\ &\text{descending coefficients }
    [1,-4,10,-17,26,-47,119,-298,\\
    &\quad 592,-878,963,-780,464,-199,59,-11,1].
\end{align*}
The final family is $h=Q_1Q_2\,Q_5^{q_E}Q_6^{q_F}$ with
\[
  q=(13.937341,12.515102,2.541409,2.068537,0.753965),\quad
  q_E=0.891271,\ q_F=0.246614,
\]
\[
  D=\max(59.198095,\ 56+0.891271\cdot12+0.246614\cdot16)=70.641076 .
\]
Admissibility is verified directly (\repofile{verify\_upper\_q6.py}, mode
\texttt{admiss}, independent \texttt{sympy}): $\deg Q_6=16>0$;
$Q_6(0)=Q_6(1)=1$; $Q_6$ irreducible/squarefree; $\gcd(Q_6,\cdot)=1$ for each of
$P_1,P_2,P_4,P_6,P_8,Q_1,Q_2$; the inter-block $\gcd(Q_6,Q_5)=1$; and the active
dictionary $W=Q_1Q_2Q_5Q_6$ (degree $84$) has $W(0)=W(1)=1$, so condition (5)
holds by unique factorization.

\subsection{The quadrature certificate}
By \eqref{eq:upperbound} it remains to enclose the integral \eqref{eq:logh} from
above. The program \repofile{verify\_upper\_q6.py} encloses
$\int_0^{2\pi}\max(A,B)\,dt$ by an adaptive outward-rounded quadrature, without
splitting the $\max$ (a $\pm\log|Q|$ split is numerically vacuous). On each cell
three guaranteed upper bounds on $\int_{\text{cell}}\max(A,B)\,dt$ are available
and the minimum is taken:
\begin{itemize}
\item \emph{flat}: $(b-a)\max(\sup A,\sup B)$---never vacuous, since a near-zero of
  $Q$ drives $B\to-\infty$ but the finite $A$ caps the $\max$;
\item \emph{straddle} (the load-bearing $O(h^2)$ bound), from
  $\max(a+x,b+y)\le\max(a,b)+\max(x,y)$ applied to the midpoint-Taylor deviation of
  each branch:
  \[
    \int_{\text{cell}}\!\max(A,B)\,dt\le
    (b-a)\max(A^\uparrow_{\mathrm{mid}},B^\uparrow_{\mathrm{mid}})
    +\max(\text{slope})\,r^2+\tfrac13\max(\text{curv})\,r^3,\ \ r=\tfrac{b-a}2;
  \]
\item \emph{midpoint} ($O(h^3)$): on cells where one branch dominates throughout
  and is tame.
\end{itemize}
Cells whose straddle deviation exceeds a cap are bisected (each bisection cuts the
$r^3$ term by $8\times$); only leaf cells contribute, with exact tiling and
all-outward rounding, so the result is a rigorous enclosure of the integral.

\begin{proof}[Proof of Theorem~\ref{thm:upper}]
The certificate run yields the frontier fully resolved ($669734$ leaves, $0$
unresolved, $6$ rounds) and
\[
  \int_0^{2\pi}G\,dt\le112.7567758427,\quad
  \int_0^1 G\,ds\le17.9457982425,\quad
  \frac{17.9457982425}{70.641076}=0.2540419719 .
\]
Therefore, by Lemma~\ref{lem:doche},
\[
  \Cee\;\le\;\log h\;\le\;0.2540419719\;<\;0.25443677=\log1.289735. \qedhere
\]
\end{proof}

\paragraph{Soundness audits}
(i) The per-cell quadrature bound dominates the true cell integral:
\texttt{selftest} checks it against a high-precision (\texttt{mpmath}, prec $160$)
reference on $200$ random cells, for both the flat and midpoint caps, with
$0/200$ violations;
(ii) anchors: $q_F=0$ recovers the previously verified family bit-identically,
$q_E=q_F=0$ the next smaller family, and $q_B=q_C=q_E=q_F=0$ Doche's base
$h=Q_1Q_2$ with $D=56$---each block is a genuine free-exponent extension;
(iii) a bogus target below the truth reports failure with the frontier resolved.

%==============================================================================
\section{Proof of Theorem~\ref{thm:main}}
%==============================================================================

Theorem~\ref{thm:lower} gives $\Cee\ge0.2524001332$ and Theorem~\ref{thm:upper}
gives $\Cee\le0.2540419719$. Combining,
\[
  0.2524001332\;\le\;\Cee\;\le\;0.2540419719,
\]
an enclosure of width $\approx1.64\times10^{-3}$.\qed

\begin{corollary}
The open interval bracketing $\Cee$ has been narrowed from
$\approx5.7\times10^{-3}$ to $\approx1.64\times10^{-3}$, a reduction of
$\approx70\%$, with both endpoints strictly improving the previously published
records.
\end{corollary}

By the strong duality of \cite[Thm.~6.5]{BMQS} the residual interval is a
truncation gap---how far the two finite optimizations have been pushed---not an
irremovable duality gap. The natural next levers are a second energy /
discriminant-moment cut (a cutting-plane iteration of the present single cut) and
a Fekete-type transfinite-diameter companion on the upper side; both lie beyond
the present single-cut certificate.

%==============================================================================
\section{Reproducibility}\label{sec:repro}
%==============================================================================

All certificates, the frozen data file, and these commands are at \repo. The
scientific Python stack suffices (\texttt{numpy}, \texttt{scipy}, \texttt{mpmath},
\texttt{sympy}); every certified inequality uses directed (outward) rounding.

\begin{center}\small
\begin{tabular}{@{}lll@{}}
\toprule
result & command (in \texttt{certificate/}) & output\\
\midrule
Thm.~\ref{thm:lower} (lower) & \texttt{python3 verify\_vec\_energy.py}
  & \texttt{CERTIFIED ... >= 0.2524001332}\\
\;\;soundness & \texttt{python3 verify\_vec\_energy.py selftest}
  & \texttt{0/620 + 0/80 violations}\\
\;\;$\Ihat$ / NSD witness & \texttt{python3 freeze\_energy.py check}
  & \texttt{Ihat <= I(mu0); top eig <= 0}\\
Thm.~\ref{thm:upper} (upper) & \texttt{python3 verify\_upper\_q6.py certify} $\ \cdots$
  & \texttt{CERTIFIED log h <= 0.2540419719}\\
\;\;admissibility & \texttt{python3 verify\_upper\_q6.py admiss}
  & all coprimality \texttt{True}\\
\;\;Flammang base & \texttt{python3 verify\_vec.py}
  & \texttt{min\_t g(t) >= 0.24874}\\
\bottomrule
\end{tabular}
\end{center}

\noindent The upper-bound command takes the exponent vector explicitly:\\
\texttt{python3 verify\_upper\_q6.py certify 13.937341 12.515102 2.541409 2.068537
0.753965 0 0 0.891271 0.246614 200000 14 1e-10}.

\medskip
\noindent\textbf{File dependencies.}
The lower-bound certificate requires \repofile{verify\_vec\_energy.py},
\repofile{verify\_vec.py}, \repofile{flammang\_table1.py},
\repofile{freeze\_energy.py} and the frozen data
\repofile{frozen\_energy.npz}. The upper-bound certificate requires
\repofile{verify\_upper\_q6.py} together with \repofile{verify\_upper\_q5.py},
\repofile{verify\_upper\_q4.py}, \repofile{verify\_upper\_q3.py},
\repofile{verify\_upper.py}, \repofile{verify\_vec.py} and
\repofile{flammang\_table1.py}. These eleven files (ten programs and the data
file) are the complete verification set; the Flammang base bound is additionally
cross-checked by the independent scalar implementation
\repofile{verify\_flammang.py}. Every other script in the repository
directory---the linear-programming search, optimizers, probes and reference
cross-checks used to \emph{find} the weights and exponents---is collected in a
\texttt{scratch/} subfolder and is not needed to \emph{verify} either theorem; no
file in the verification set imports any of them.

\medskip
\noindent\textbf{Live reproduction note.}
Theorem~\ref{thm:lower}'s certificate was re-run during preparation and reproduced
$0.2524001332$ exactly (4.1\,s); the soundness selftest reproduced
$0/620+0/80$; \texttt{freeze\_energy.py check} reproduced
$\Ihat=-0.2111616260\le I(\mu_0)_{\mathrm{ref}}=-0.2111600446$ with NSD top
eigenvalue $+4.0\times10^{-16}$; the upper admissibility check reproduced all
coprimality flags \texttt{True}. Theorem~\ref{thm:upper}'s full quadrature
reproduces $0.2540419719$.

%==============================================================================
\begin{thebibliography}{99}

\bibitem{BMQS} J.~Burgos Gil, R.~Menares, B.~Qu, M.~Sombra,
  \emph{Closing the gap around the essential minimum of height functions with
  linear programming}, preprint, arXiv:2601.18978 (2026).

\bibitem{Doc01a} Ch.~Doche,
  \emph{On the spectrum of the Zhang--Zagier height},
  Math.\ Comp.\ \textbf{70} (2001), no.~233, 419--430.

\bibitem{Doc01b} Ch.~Doche,
  \emph{Zhang--Zagier heights of perturbed polynomials},
  J.\ Th\'eor.\ Nombres Bordeaux \textbf{13} (2001), no.~1, 103--110.

\bibitem{F18} V.~Flammang,
  \emph{On the Zhang--Zagier measure},
  Int.\ J.\ Number Theory \textbf{14} (2018), no.~10, 2663--2671.

\bibitem{OSS} J.~Orloski, N.~T.~Sardari, A.~Smith,
  \emph{On the logarithmic energy and the Schur--Siegel--Smyth trace problem},
  preprint, arXiv:2401.03252.

\bibitem{OT23} J.~Orloski, F.~Talebizadeh Sardari,
  \emph{Multivariate integrality and Galois orbits of algebraic integers},
  preprint (2023).

\bibitem{S80} C.~J.~Smyth,
  \emph{On the measure of totally real algebraic integers},
  J.\ Austral.\ Math.\ Soc.\ Ser.\ A \textbf{30} (1980/81), no.~2, 137--149.

\bibitem{Zag93} D.~Zagier,
  \emph{Algebraic numbers close to both $0$ and $1$},
  Math.\ Comp.\ \textbf{61} (1993), no.~203, 485--491.

\bibitem{Zha95} S.~Zhang,
  \emph{Positive line bundles on arithmetic varieties},
  J.\ Amer.\ Math.\ Soc.\ \textbf{8} (1995), no.~1, 187--221.

\end{thebibliography}

\end{document}
